\documentclass[11pt]{letter}
\usepackage[margin=1in]{geometry}
\usepackage[T1]{fontenc}
\usepackage[utf8]{inputenc}
\usepackage{parskip}
\usepackage[hidelinks]{hyperref}

\signature{Quang-Vinh Dang}
\address{Quang-Vinh Dang\\
British University Vietnam\\
Hung Yen, Vietnam\\
\href{mailto:vinh.dq4@buv.edu.vn}{vinh.dq4@buv.edu.vn}}

\begin{document}

\begin{letter}{The Editor-in-Chief\\ \emph{Borsa Istanbul Review}}

\opening{Dear Editor-in-Chief,}

I am pleased to submit for your consideration the manuscript entitled
\textbf{``What Financial Statements Can and Cannot Reveal About Monetary
Transmission to Corporate Balance Sheets: Evidence from Vietnam and
T\"{u}rkiye''} for publication in \emph{Borsa Istanbul Review}.

\textbf{Motivation.} A productive literature shows that firms holding
floating-rate debt bear the brunt of monetary tightening. It rests almost
entirely on supervisory credit registries that record contract-level repricing
terms. Most emerging markets have no such registry available to researchers;
what they have is published financial statements. This paper asks, formally and
empirically, whether the accounting record can substitute---and it matters,
because the answer determines how far this literature can travel.

The paper takes as its point of departure Şengül and Çinko's recent
contribution to this journal (2026, 26(4), 100837), which uses Turkish
administrative data and a Double Machine Learning design to identify
heterogeneous effects of the June 2023 tightening. My aim is complementary
rather than critical: I ask what becomes possible, and what does not, when the
administrative layer they exploit is unavailable. I use their setting as one of
two countries and adopt their substantive question as the target.

\textbf{What the paper does.} I develop a measurement framework in which
exposure is estimated rather than observed, deriving seven propositions, and
then take their predictions to harmonised panels of 693 Vietnamese and 585
Turkish listed non-financial firms over 2010--2025. The two countries provide a
natural contrast in dose: T\"{u}rkiye tightened from 9\% to 47.5\% while Vietnam
tightened mildly and reversed within a year.

The contributions are of three kinds:

\begin{enumerate}
\item \textbf{A validated price measure.} A simple implied borrowing cost
recovers each central bank's policy path without being calibrated to it,
peaking in T\"{u}rkiye in 2024 at 31.2\% as policy peaks at 47.5\%. I use it to
estimate pass-through of roughly 35\% alongside a halving of median interest
coverage.

\item \textbf{Two data hazards of direct relevance to this journal's
readership.} First, a widely used Turkish vendor field bundles foreign-exchange
revaluation into interest expense, so measures built on it peak in the lira
crises of 2018 and 2021 rather than at the policy peak---a researcher would
recover currency shocks and report them as monetary transmission. Second, and
more consequentially, T\"{u}rkiye's inflation-accounting restatement from 2022
indexes non-monetary assets while leaving monetary debt nominal; because the
factor depends on asset vintage, the distortion enters a
difference-in-differences as an exposure-by-post interaction that fixed effects
cannot absorb. I show this generates highly significant, wrongly signed
estimates. This threatens a broad class of studies using Turkish firm-level
balance-sheet data spanning 2022.

\item \textbf{A transferable diagnostic.} I show that the split-half correlation
of independently estimated exposures \emph{is} the attenuation factor, not
merely a diagnostic correlated with it, and that decomposing the variance
separates imprecision from impermanence. In my data, conventional standard
errors implied reliabilities of 0.78--0.99 where the permanent component was in
fact zero. A power calculation would have endorsed the design.
\end{enumerate}

\textbf{On the central result.} The paper reports that no firm-level exposure
gradient survives once contaminated outcomes are removed and the pre-period is
long enough to test parallel trends. I recognise that null findings warrant
scrutiny, and the manuscript addresses this directly rather than by assertion:
three significant results appeared during the analysis and each is shown to be
an artefact---small-sample noise, accounting restatement, and pre-existing
trend---by diagnostics a conventional specification would not have run. The
framework explains the null structurally rather than attributing it to power.

\textbf{Fit with the journal.} The manuscript is an empirical study of corporate
finance and monetary economics in emerging markets, comparative in design, and
its two data-hazard findings bear directly on work using Turkish firm-level
data. It engages a paper this journal has recently published and extends its
question to settings without registry access.

\textbf{Declarations.} This manuscript is original, has not been published
previously, and is not under consideration for publication elsewhere. There are
no conflicts of interest to declare and the research received no specific
funding. I am the sole author. In accordance with the journal's double
anonymized review policy, identifying information is confined to the separate
title page and the manuscript file contains none.

Thank you for considering this submission.

\closing{Yours sincerely,}

\end{letter}
\end{document}
